% Part: lambda-calculus
% Chapter: introduction
% Section: term-revisited

\documentclass[../../../include/open-logic-section]{subfiles}

\begin{document}

\olfileid{lam}{syn}{tr}
\olsection{ترمونه د $\alpha$-معادلتوب د ټولګيو په توګه}

له دې وروسته به ترمونه د $\alpha$-معادلتوب تر کچې څېړو۔ يعنې کله
چې يو ترم ليکو، د هغه د $\alpha$-معادلتوب ټولګے مو مراد وي۔
د بېلګې په توګه، له $\lambd[a][\lambd[b][a c]]$ څخه د ټولو هغو
ترمونو سټ مرادوو چې ورسره $\alpha$-معادل وي؛ په دې کښې
$\lambd[a][\lambd[b][a
    c]]$ او $\lambd[b][\lambd[a][b c]]$ هم راځي، او داسې نور۔

همدارنګه، په پخوانيو برخو کښې د $N, Q$ په څېر توري د يوه ترم
دپاره کارېدل؛ له دې وروسته به هر يو د يوه ټولګي نښه وي۔
راتلونکې څېړنه به د ترمونو پر ځاے د دغو ټولګيو په باب وي۔
د $x, y$ په څېر توري بيا هم متغيرونه ښيي۔

د~$M$ د ټولګي د يوه خپل‌سري !!{element} د ښودلو لپاره
$\rep{M}$ ليکو۔ که له يوه زياتو ته اړتيا وي،
$\rep{M}[0], \rep{M}[1]$ او داسې نور ليکو۔
\textbf{د ژباړې د نسخې سمون (OLLAM-027):} سرچينې د «او داسې
نور» عبارت د رياضي په چاپېريال کښې راوړے و؛ دلته دا عادي
عبارت په پښتو متن کښې له رياضي بهر ليکل شوے دے۔

د وينا د اسانولو دپاره د ترمونو پخوانۍ نښې بيا کاروو۔ پر
ټولګيو لاندې تعريف لرو:
\begin{defn}
  \begin{enumerate}
  \item $\lambd[x][N]$ هغه ټولګے دے چې
    $\lambd[x][\rep{N}]$ پکې شامل وي۔
  \item $PQ$ هغه ټولګے دے چې $\rep{P}\rep{Q}$ پکې شامل وي۔
  \end{enumerate}
\end{defn}

دا تعريفونه د استازي له ټاکنې خپلواک دي؛ ځکه چې
$\alpha$-بدلون له دغو جوړښتونو سره سازګار دے۔

\begin{defn} \ollabel{def:fv}
  د $\alpha$-معادلتوب د ټولګي~$M$ \emph{ازاد متغيرونه}، يا
  $FV(M)$، داسې تعريفېږي: $FV(\rep{M})$۔
\end{defn}

دا تعريف د استازي له ټاکنې خپلواک دے، ځکه
$FV(\rep{M}[0]) =
FV(\rep{M}[1])$؛ دا په \olref[alp]{thm:fv} کښې ښودل شوي دي۔

پر ټولګيو د تعويض دپاره هم پخوانۍ نښه کاروو:
\begin{defn} \ollabel{def:sub}
  په $M$ کښې د $y$ پر ځاے د $R$ \emph{تعويض}، يعنې
  $\Subst{M}{R}{y}$، د داسې هر $\rep{M}$ او $\rep{R}$
  دپاره چې تعويض ئې ټاکلے وي، په
  $\Subst{\rep{M}}{\rep{R}}{y}$ تعريفېږي۔
  \textbf{د ژباړې د نسخې سپيناوی (OLLAM-028):} په دې تعريف
  کښې د وروستي خام ترم د الفا-معادلتوب ټولګے د پايلې
  په توګه اخلو؛ سرچينه د پايلې دغه ټولګي‌والي په څرګنده
  نۀ وايي۔
\end{defn}

سرچينه د دې تعريف د استازي له ټاکنې خپلواکي
\olref[alp]{cor:sub} ته راجع کوي۔
\textbf{د ژباړې د نسخې يادونه (OLLAM-029):} د هغې پايلې
تر شا د يکتا والي د ثبوت تشه په مخکنۍ برخه کښې څرګنده
شوې ده؛ نو دلته دا حواله د بشپړ شوي ثبوت په توګه نۀ اخلو۔

وګورئ چې دا تعريف زمونږ استدلال څومره اسانوي۔ د بېلګې په توګه:
\begin{align}
  \Subst{\lambd[x][x]}{y}{x} & \ollabel{eq:1}\\
  &= \Subst{\lambd[z][z]}{y}{x} \ollabel{eq:2}\\
  &= \lambd[z][\Subst{z}{y}{x}] \\
  &= \lambd[z][z]
\end{align}
\textbf{د ژباړې د نسخې سمون (OLLAM-030):} د سرچينې د لومړۍ
کرښې په پاے کښې د مساوات نښه بې ښي اړخه پاتې وه۔ هغه
نښه لرې شوه؛ نورې کرښې او د حوالو نښې هماغسې دي۔

که دا لا هم پر خامو ترمونو تعويض وګڼو، نو \olref{eq:1}
ناټاکلے دے۔ خو لکه پورته چې وويل شول، اوس پر ټولګيو
تعويض څېړو۔ له همدې امله \olref{eq:2} روا دے:
$\lambd[x][x]$ په $\lambd[z][z]$ بدلولے شو، ځکه چې
دواړه په يوه ټولګي کښې دي۔

له همدې امله به له دې وروسته فرض کوو چې غوره شوي
استازي د تعويض اړين شرطونه پوره کوي۔ مثلاً کله چې
$\Subst{\lambd[x][N]}{R}{y}$ ووينو، فرض کوو چې
استازے $\lambd[x][N]$ داسې غوره شوے دے چې
$x \neq y$ او $x \notin FV(R)$ وي۔

څرنګه چې $\lambd[x][x]$ ته «ټولګے» ويل لږ نااشنا ښکاري،
له دې وروسته به دې ټولګيو ته $\Lambda$-ترمونه وايو
(يا د دې برخې په پاتې متن کښې يوازې «ترمونه»)، څو
له هغو $\lambda$-ترمونو سره ئې بېل وساتو چې له پخوا
ورسره اشنا يو۔

\begin{editorial}
  لا له خامو ترمونو سره مخه نۀ شو ښۀ کولے: د
  $\Lambda$-ترمونو ټول تعريف پر $\lambda$-ترمونو ولاړ دے،
  او تر اوسه مو پر $\Lambda$-ترمونو د تابعګانو د تعريف
  طريقه نۀ ده ورکړې۔ له همدې امله دا ډول تابعګانې
  لومړے پر $\lambda$-ترمونو تعريفوو او بيا ئې
  $\Lambda$-ترمونو ته «لېږدوو»، لکه د تعويض په حالت کښې۔
  خو فرض کوو چې لوستونکے په بديهي ډول پوهېږي چې پر
  $\Lambda$-ترمونو تابعګانې څنګه تعريفېدے شي۔
  \textbf{د ژباړې د نسخې سپيناوی (OLLAM-031):} داسې
  لېږد يوازې هغه وخت پر ټولګيو ښه تعريف لري چې پايله
  د استازي له ټاکنې خپلواکه وي؛ سرچينه دلته دا اړين
  شرط نۀ بيانوي۔
\end{editorial}
\end{document}
